% 04-Dataset.tex -- data construction and experimental setup.
%
% AAAI COMPLIANCE: no added packages, and none of the commands aaai2027
% forbids (no \setlength, \vspace, \vskip, \newpage, \tiny, ...). Column
% padding is trimmed with @{} rather than by resetting \tabcolsep.
%
% SCOPE: experimental detail only -- what was run, on what data, scored
% how. No findings, no interpretation; those live in the Results section.
%
% \text comes from amsmath, which this template does not load; the
% fallback keeps the file self-contained if any \text survives an edit.
% It is a no-op when amsmath is present.
\providecommand{\text}[1]{\mathrm{#1}}

\section{Data}

\subsection{Universe and Source}

All data are daily CRSP records for U.S.\ common equities. From the
mid-- and upper-mid-capitalisation segment we form four disjoint panels,
denoted $V1$--$V4$, of 50 tickers each; the four share no constituents,
so 200 distinct names are used in total. Each panel is stored as one CSV
per ticker plus a panel-level date index carrying the aligned raw price
and transaction-cost columns for every name.

Every panel ends on 2024-12-31 but they begin on different dates,
because a panel starts only once all 50 of its constituents have priced
history: $V1$ runs from 2004-08-16, $V2$ from 2004-12-17, $V3$ from
2004-10-27 and $V4$ from 2004-02-09, giving 5{,}130, 5{,}043, 5{,}079 and
5{,}257 trading days respectively. The differing lengths affect only the
burn-in each panel receives; the tuning and trading periods are the same
calendar spans for all four, so no arm gains or loses scored days.

Panel membership is fixed for the whole sample: a ticker is either in a
panel for its entire history or not in it at all. This removes
index-reconstitution effects from the comparison, at the cost of
conditioning the universe on names that survived to the end of the
sample. All methods trade the identical universe, so the conditioning
applies equally to every arm reported.

\subsection{Features and Prediction Target}

Each model sees seven per-stock features, all computed causally from
information available at or before the formation day:

\begin{itemize}\itemsep2pt
\item \texttt{RET} --- the raw daily return.
\item \texttt{RET\_CS\_IN} --- the cross-sectional rank-normal score of
      \texttt{RET} on the same day, i.e.\ $\Phi^{-1}\!\big((\mathrm{rank}_t(i)-0.5)/N\big)$
      over the $N$ names priced that day.
\item \texttt{BA\_SPREAD} --- the quoted bid--ask spread, the gap between
      the best buying and selling price, which proxies how costly the
      stock is to trade.
\item \texttt{ILLIQUIDITY} --- Amihud illiquidity: absolute return over
      dollar volume, smoothed over 21 days and logged. It is large when
      a small amount of trading moves the price a lot.
\item \texttt{REV21\_1} --- short-horizon reversal: the compounded return
      over the 20 days before the formation day, excluding it. Recent
      winners tend to give some of it back, so this is predictive with a
      negative sign.
\item \texttt{TURN\_RATIO} --- turnover acceleration: recent trading
      activity relative to its own longer-run level, as the mean of
      $\log(1+\mathrm{turnover})$ over the last 5 days minus the same
      mean over the last 63. It rises when a stock starts attracting
      unusual attention.
\item \texttt{VOL21} --- realised volatility, the standard deviation of
      \texttt{RET} over the trailing 21 days.
\end{itemize}

These are standard predictors in the empirical asset pricing literature
rather than novel ones; the contribution of this paper is in the
architecture search, so the feature set is deliberately conventional.

The last five of these are additionally mapped, per date, to their
cross-sectional rank on $[-1,1]$, so that no feature carries a level that
drifts across the two decades of the sample.

The prediction target is \texttt{RET\_CS}: the next day's return, rank-
normalised across the cross-section of that day by the same $\Phi^{-1}$
transform. Predicting a cross-sectional rank rather than a raw return
makes the target stationary by construction and matches how the
predictions are consumed --- every trading rule in this paper reads only
the ordering of the predicted cross-section.

Realised returns used for scoring are CRSP \texttt{RET}, which is adjusted
for splits and dividends. Prices and per-name transaction costs enter only
through the cost model below.

\section{Experimental Setup}

\subsection{Online Evaluation Protocol}

Every method is evaluated prequentially: on each scored day the model
predicts the next day's cross-section, the realised outcome is then
revealed, and the model updates. No model ever sees a label before
predicting it, and no model is refit on the scored span with hindsight.

ONE-NAS advances on a weekly clock. One generation corresponds to a
\texttt{window\_step} of 5 trading days; each genome is trained on
sequences of length $L=40$ and validated on the 5 most recent windows.
The number of training windows is set per panel so that the clock lands on
the intended calendar span --- 712, 695, 702 and 738 windows for
$V1$--$V4$ on the evaluation span, and 511, 493, 501 and 536 on the
tuning span.

The sample is divided into three disjoint spans, used for three distinct
purposes:

\begin{itemize}\itemsep2pt
\item \textbf{Burn-in}, through 2021-12-31. The population evolves online
      throughout, but nothing is traded or scored. A further 50 warm-up
      generations precede any recorded prediction, so the population that
      begins trading is mature rather than cold-started.
\item \textbf{Tuning}, 2016-01-01 to 2019-12-31. Every configuration
      decision --- hyperparameter screen, island count --- is made here,
      disjoint from the trading period.
\item \textbf{Trading}, 2022-01-01 to 2024-12-31. All reported results
      are on this period, which is also the period the prior study on
      these panels evaluated, so the two are directly comparable.
\end{itemize}

\subsection{Search Configuration and Prediction Rule}

The ONE-NAS settings are listed in Table~\ref{tab:hps} and were selected
on the tuning span. Node types available to the search are simple,
UGRNN, MGU, GRU, $\Delta$-RNN and LSTM recurrent cells.

A single run yields, at every generation, both the elite population of
each island and the single best genome overall, so the prediction rule is
a scoring-time choice rather than a property of the run. Two rules are
reported. The \emph{single champion} rule takes the global-best genome, as
in prior ONE-NAS work. The \emph{ensemble} rule takes one champion per
island --- the lowest-validation-error elite of each --- converts each
member's predicted cross-section to ranks, and averages the ranks. Rank
averaging is used rather than averaging predictions because the members
are separately trained networks whose output scales are not comparable.

\subsection{Hyperparameter Search}

Every configuration decision was made on the tuning span and is listed
here in full, so that the selected values in Table~\ref{tab:hps} can be
read against what was actually considered.

The search is single-axis around the registered control configuration:
each cell changes one setting and is compared against the control on
identical (panel, seed) cells. Sixteen configurations were trained ---
the control plus fifteen variants:

\begin{itemize}\itemsep2pt
\item \textbf{Fitness signal.} Selection on validation MSE (control),
      on cross-sectional rank IC, and on a gated rank IC.
\item \textbf{Prediction target.} 1-day cross-sectional rank-normal
      return (control) and the 5-day equivalent, the latter also in
      combination with gated-IC selection.
\item \textbf{Population structure.} Elite/generated ratios of 8/5
      (control), 5/10 and 5/15; repopulation every 25, 50 (control) and
      100 generations.
\item \textbf{Training budget.} 2000 (control) and 600 training
      subsequences per genome; 10 (control) and 20 backpropagation
      epochs per genome; one (control) and two generations per window.
\item \textbf{Replay sampling.} Prioritised replay at
      $\alpha \in \{0.4, 0.6\ (\text{control}), 0.8\}$, at
      $\lambda \in \{0.007\ (\text{control}), 0.021\}$, and uniform
      sampling with prioritisation disabled.
\item \textbf{Island count.} 8, 16, 20, 40, 50 and 60
      (Table~\ref{tab:tunewidth}).
\end{itemize}

Selection used the tuning-span rank IC as the primary objective with the
registered book's net return and Sharpe as an economic co-primary, and
required a paired improvement over the control to adopt a change. No
single-axis variant met that bar, so the control configuration is the one
reported; the island count is the one axis where a change was adopted,
and the value chosen sits on the plateau of Table~\ref{tab:tunewidth}
rather than at its nominal maximum. Table~\ref{tab:invariance} reports
every trained cell.

Several axes were deliberately not searched, having been settled by
earlier work on this data: sequence length, the number of validation
sets, the node-type library, window geometry, and the ensemble's width
and combination rule.

\subsection{Baselines}

Four online baselines trade the same books on the same panels under the
same costs. All are single networks; the reported cell is the mean over
independent runs, not the return of an ensemble of those runs.

\begin{itemize}\itemsep2pt
\item \textbf{Online LSTM} --- one LSTM, 8 hidden units, sequence length
      10, Adam at $10^{-3}$, one gradient step per trading day with a
      1000-day replay buffer.
\item \textbf{Online GRU} --- as above with a GRU cell and sequence
      length 5.
\item \textbf{Periodic LSTM} --- the same LSTM architecture refit from
      scratch each month for up to 4000 steps, rather than updated daily.
\item \textbf{Online AR} --- a first-order autoregression fitted by online
      gradient descent at $\eta=10^{-4}$. Deterministic, so it contributes
      one run per panel rather than one per seed.
\end{itemize}

All baselines predict the same cross-sectionally demeaned target as
ONE-NAS.

\subsection{Trading Books}

A \emph{book} is the rule that turns a day's predicted ranking into
actual positions, and it is what converts prediction quality into money.
All three books used here are \emph{long--short}: they buy (go
\emph{long}) the names predicted to do best and sell borrowed shares of
(go \emph{short}) the names predicted to do worst, in equal dollar
amounts. Equal amounts on each side make the book roughly
\emph{dollar-neutral}, so it profits from the top names outperforming the
bottom ones rather than from the market rising, and a rising or falling
market largely cancels out. This is why the books can be compared against
each other on skill rather than on market exposure, and why their returns
are not comparable to simply holding the market.

One registered primary construction is used, with two further
constructions reported as sensitivities.

The \textbf{sleeves} book is the registered primary: Jegadeesh--Titman
overlapping portfolios. Each day a new \emph{sleeve} --- one small
long--short portfolio --- goes long the top $k=10$ names and short the
bottom 10 by predicted rank, at $(\$100/H)/k$ per name per side with
holding period $H=10$ days; the sleeve formed $H$ days earlier is
retired. Ten sleeves are therefore live at once, each holding for ten
days, so the book turns over about a tenth of itself per day rather than
rebuilding daily. Costs are charged on the change in \emph{aggregate}
per-name position, so a name held by several sleeves at once is netted
rather than bought and sold repeatedly. Because only the ordering of each
day's predictions is read, every metric is invariant to any strictly
monotone transform of a model's outputs --- a model that doubles all its
predictions, or applies any order-preserving rescaling, trades
identically.

The \textbf{banded} book enters a name at predicted rank $\le 10$ and
holds it until its rank passes 35, following the buy/hold-spread
construction used to reduce turnover in the transaction-cost literature.
The \textbf{Algorithm 1} book is the conditional long--short rule of the
prior study, which rebalances only when all ten top predictions are
positive and all ten bottom predictions negative.

Books are funded with \$100 of capital and, when fully invested, hold
\$100 long and \$100 short, so \$200 of stock changes hands in total ---
the \emph{gross notional}.

Trading is not free, and the cost is charged explicitly. Every trade pays
the \emph{bid--ask spread}: the gap between what a buyer pays and a
seller receives, which is wide for thinly traded stocks and narrow for
liquid ones. We take that cost per name from the CRSP cost field and
apply it to the value actually traded, rather than assuming one flat rate
for every stock, so illiquid names are correctly penalised. All returns
reported in this paper are \emph{net}: costs have already been deducted.

\subsection{Metrics}

Two families are reported: how good the predictions are, and what the
resulting book earns.

\emph{Prediction quality} is the \textbf{rank IC} (information
coefficient): on each day, the Spearman rank correlation between the
predicted ordering of the cross-section and the realised one, averaged
over scored days. It is the natural accuracy measure here because the
books read only the ordering. Values are small by construction --- daily
cross-sectional return is mostly noise, and a rank IC of $0.01$--$0.02$
is a usable signal in this literature, not a weak one.

\emph{Economics} covers four quantities:

\begin{itemize}\itemsep2pt
\item \textbf{Net return}, the cumulative profit and loss over the period
      as a percentage of the starting \$100, after costs.
\item \textbf{Sharpe ratio}, the mean daily return divided by its
      standard deviation, annualised. It measures return per unit of
      volatility, so a strategy that earns the same amount more smoothly
      scores higher. Roughly, above 1 is good for a strategy of this
      kind.
\item \textbf{Maximum drawdown} (MDD), the largest peak-to-trough fall in
      cumulative profit over the period --- the worst loss an investor
      who entered at the worst moment would have sat through.
\item \textbf{Turnover}, mean daily traded value as a fraction of gross
      notional. A turnover of $0.12$ means about $12\%$ of the book is
      replaced each day; lower turnover means less exposure to
      transaction costs.
\end{itemize}

Finally, a return can look impressive merely by taking well-known risks,
so we check whether it survives that. Each arm's daily book return is
regressed on the four standard \emph{factors} of the empirical asset
pricing literature --- the market, firm size, value, and momentum
(Fama--French three-factor plus momentum). The slope on the market factor
is the \textbf{beta}, the portfolio's sensitivity to the market as a
whole: a beta near $1$ moves with the market, and one near $0$ is
insulated from it. The intercept is the \textbf{alpha}, the average
return left over once those four exposures are accounted for, and it is
alpha rather than raw return that indicates genuine skill. Standard
errors use the Newey--West correction at lag 10, which allows for daily
returns being autocorrelated rather than independent.

\subsection{Replication and Statistical Comparison}

The unit of replication is one (panel, seed) cell. Unless stated
otherwise, ONE-NAS arms comprise 10 seeds $\times$ 4 panels $=40$ cells.
The baselines were run at higher replication --- 48 seeds for online LSTM,
40 for online GRU and periodic LSTM --- and their own means are reported
at that full replication.

Comparisons between arms are paired within (panel, seed) cells, so that
panel and seed effects cancel, and the reported $t$ is a paired $t$ over
those cells. Where two arms do not share a seed set, the pairing is
restricted to the cells they have in common and $n$ is stated.

\subsection{Compute Environment}

All searches ran on ACCESS Anvil. Each run occupied one node of the
\texttt{wholenode} partition with 128 MPI ranks, one acting as the master
that maintains the island populations and distributes genomes, and the
remaining 127 as workers that train and evaluate them. The island count
is a search parameter and is independent of the rank count, so the worker
pool is identical across every configuration reported. Baselines are
single-process and were run on the same hardware.

\subsection{Pre-Registration}

The scored configuration, evaluation span, trading book and metrics were
committed to the repository before the evaluation span was scored, and
each subsequent look at the data was declared in advance with its
decision rule. Configuration changes require a dated amendment written
before the changed configuration is scored. The registered protocol,
the hyperparameter screen with its per-cell outcomes, and every declared
exploratory look are included in the code release.
